\section{Wyniki operacji CRUD}

Macierz obejmuje \SizesList{} biletów i~dała \ScenarioRuns{} pomiarów
scenariuszy oraz \ContentionRuns{} rund testu współbieżności. Wszystkie czasy
w~tym rozdziale są czasami po stronie serwisu bazy, w~milisekundach, uśrednionymi
z~5 powtórzeń.

\subsection{Import danych i rozmiar}

\begin{table}[H]
\centering
\caption{Czas importu i~tworzenia indeksów drugorzędnych [s].}
\label{tab:load}
\small
\tbinput{generated/tables/load.tex}
\end{table}

\begin{table}[H]
\centering
\caption{Zajętość danych [MB]. Dla baz dyskowych jest to rozmiar plików bazy,
dla Valkey zużycie pamięci (\texttt{used\_memory}).}
\label{tab:storage}
\small
\tbinput{generated/tables/storage.tex}
\end{table}

Import 10 mln biletów najszybciej wykonuje Valkey
(\lval{valkey}{10000000}{load}\,s), potem PostgreSQL
(\lval{postgres}{10000000}{load}\,s, binarny \texttt{COPY}), MariaDB
(\lval{mariadb}{10000000}{load}\,s, wielowierszowe \texttt{INSERT}) i~MongoDB
(\lval{mongo}{10000000}{load}\,s). MongoDB płaci za \emph{write concern}
\texttt{majority} z~dziennikiem i~za budowę dokumentów BSON z~osadzonymi
tablicami. Czas importu rośnie z~rozmiarem w~przybliżeniu liniowo we
wszystkich bazach.

Tworzenie 15 indeksów na 10 mln biletów trwa w~PostgreSQL
\lval{postgres}{10000000}{indexes_on}\,s, w~MariaDB
\lval{mariadb}{10000000}{indexes_on}\,s, a~w~MongoDB
\lval{mongo}{10000000}{indexes_on}\,s. W~Valkey ,,indeksy'' buduje aplikacja:
musi przeczytać przez \texttt{SCAN} wszystkie 1,67 mln rezerwacji i~zapisać
zbiory posortowane, co zajmuje \lval{valkey}{10000000}{indexes_on}\,s.

Najmniej miejsca zajmuje MongoDB dzięki kompresji WiredTiger (Snappy):
\sval{mongo}{10000000}{off}\,MB bez indeksów wobec
\sval{postgres}{10000000}{off}\,MB w~PostgreSQL i~\sval{mariadb}{10000000}{off}\,MB
w~MariaDB. Valkey przy 10 mln biletów zajmuje \sval{valkey}{10000000}{off}\,MB
pamięci RAM. Cały zbiór musi się w~niej zmieścić, co jest praktyczną granicą
tego modelu. Wartości ,,z~indeksami'' obejmują też wiersze dopisane przez
scenariusze z~pierwszego przebiegu (np. C3 importuje \num{10000} biletów
w~każdym powtórzeniu).


\subsection{Czasy scenariuszy}

Tabele~\ref{tab:s500}--\ref{tab:s10m} zawierają średni czas każdego
scenariusza. Pogrubiona jest najlepsza wartość w~danym wariancie indeksów.
,,t/o'' oznacza przekroczenie limitu 60~s.

\newcommand{\scentable}[3]{%
\begin{table}[htbp]
\centering
\caption{Średni czas scenariuszy [ms], #2 biletów.}
\label{#3}
\resizebox{\linewidth}{!}{\tbinput{generated/tables/scenarios_#1.tex}}
\end{table}}
\scentable{500000}{500 tys.}{tab:s500}
\scentable{1000000}{1 mln}{tab:s1m}
\scentable{10000000}{10 mln}{tab:s10m}

\newcommand{\scenplot}[2]{%
  \nextgroupplot[title={#1 #2}]
  \addplot[pg] table[x=tickets, y=postgres] {generated/plots/scen_#1_#2.dat};
  \addplot[maria] table[x=tickets, y=mariadb] {generated/plots/scen_#1_#2.dat};
  \addplot[mongo] table[x=tickets, y=mongo] {generated/plots/scen_#1_#2.dat};
  \addplot[valkey] table[x=tickets, y=valkey] {generated/plots/scen_#1_#2.dat};}

\begin{figure}[p]
\centering
\begin{tikzpicture}
\begin{groupplot}[group style={group size=3 by 4, horizontal sep=1.3cm, vertical sep=1.5cm},
  width=0.36\linewidth, height=4.4cm, xmode=log, ymode=log, grid=major, grid style={gray!25},
  tick label style={font=\scriptsize}, title style={font=\footnotesize}, unbounded coords=jump, log ticks with fixed point, every axis plot/.append style={mark options={scale=0.7}},
  xtick={500000,1000000,10000000}, xticklabels={0{,}5,1,10}]
  \scenplot{R1}{off}
  \scenplot{R2}{off}
  \scenplot{R4}{off}
  \scenplot{R1}{on}
  \scenplot{R2}{on}
  \scenplot{R4}{on}
  \scenplot{R6}{off}
  \scenplot{U6}{off}
  \scenplot{D4}{off}
  \scenplot{R6}{on}
  \scenplot{U6}{on}
  \scenplot{D4}{on}
\end{groupplot}
\end{tikzpicture}

\vspace{2mm}
\footnotesize
\textcolor{cpg}{$\bullet$}~PostgreSQL\quad
\textcolor{cmaria}{$\blacksquare$}~MariaDB\quad
\textcolor{cmongo}{$\blacktriangle$}~MongoDB\quad
\textcolor{cvalkey}{$\blacklozenge$}~Valkey
\caption{Czas wybranych scenariuszy [ms] w~funkcji liczby biletów w~mln (obie osie
logarytmiczne). ,,off'': bez indeksów drugorzędnych, ,,on'': z~indeksami.}
\label{fig:scen}
\end{figure}

\paragraph{Operacje punktowe.} Wstawienie i~zmiana pojedynczego rekordu
po kluczu (C1, C2, C4--C6, U1, U2, U4, U5, J3) we wszystkich bazach trwają
poniżej 2\,ms i~nie zależą od rozmiaru zbioru. Pojedyncza rezerwacja (C1) na
10 mln biletów trwa \val{valkey}{10000000}{off}{C1}\,ms w~Valkey,
\val{postgres}{10000000}{off}{C1}\,ms w~PostgreSQL,
\val{mariadb}{10000000}{off}{C1}\,ms w~MariaDB i~\val{mongo}{10000000}{off}{C1}\,ms
w~MongoDB. MongoDB jest tu najwolniejszy, bo każdy zapis czeka na
potwierdzenie z~dziennika (\texttt{majority}), a~C1 dodatkowo jest
transakcją wielodokumentową. Valkey nie zapisuje niczego na dysk, więc jego
przewaga w~zapisach wynika częściowo z~innego poziomu trwałości.

\paragraph{Zapytania bez indeksów rosną liniowo.} W~bazach relacyjnych
i~w~MongoDB każde zapytanie po kolumnie innej niż klucz główny wymaga pełnego
skanu. Wolne miejsca na wydarzenie (R2) w~PostgreSQL rosną z~\val{postgres}{500000}{off}{R2}\,ms
(500 tys.) do \val{postgres}{10000000}{off}{R2}\,ms (10 mln). W~MariaDB rosną
z~\val{mariadb}{500000}{off}{R2} do \val{mariadb}{10000000}{off}{R2}\,ms,
a~w~MongoDB z~\val{mongo}{500000}{off}{R2} do \val{mongo}{10000000}{off}{R2}\,ms.
W~Valkey ten sam scenariusz trwa stale około \val{valkey}{10000000}{off}{R2}\,ms,
bo zbiór wolnych miejsc wydarzenia jest podstawową strukturą modelu, a~nie
indeksem.

\paragraph{MariaDB bez indeksów jest najwolniejsza z~relacyjnych.} Przy 10 mln
biletów MariaDB wykonuje R1 w~\val{mariadb}{10000000}{off}{R1}\,ms, a~PostgreSQL
w~\val{postgres}{10000000}{off}{R1}\,ms. Złączenie rezerwacji z~biletami bez
indeksu MariaDB realizuje pętlą zagnieżdżoną z~buforem bloków (\texttt{ALL} na
tabeli \texttt{bilety}, ok. 10 mln przeczytanych wierszy, tab.~\ref{tab:plans-rel}).
PostgreSQL używa złączenia haszowego i~równoległego skanu. Ta sama różnica
widać w~R5, U3, U6 i~D1--D6.

\paragraph{Osadzanie zamiast złączeń.} Odczyt rezerwacji z~biletami (R1)
w~MongoDB trwa \val{mongo}{10000000}{off}{R1}\,ms, a~w~Valkey
\val{valkey}{10000000}{off}{R1}\,ms, także bez indeksów drugorzędnych.
Bilety są osadzone w~dokumencie albo w~JSON-ie rezerwacji, więc wystarcza
wyszukiwanie po kluczu głównym. Ten sam efekt widać przy anulowaniu
rezerwacji (D1): MongoDB zwalnia bilety po identyfikatorach z~osadzonej
migawki, przez indeks \texttt{\_id}, w~\val{mongo}{10000000}{off}{D1}\,ms.
PostgreSQL bez indeksu na \texttt{bilety.rezerwacja\_id} potrzebuje do tego
\val{postgres}{10000000}{off}{D1}\,ms, a~MariaDB \val{mariadb}{10000000}{off}{D1}\,ms.
To przykład kompromisu normalizacji i~denormalizacji: odczyt agregatu jest
tani, ale dane biletów są przechowywane w~dwóch miejscach.

\paragraph{Agregacje.} Ranking 10 najpopularniejszych wydarzeń (R6) wymaga
w~bazach relacyjnych i~w~MongoDB przejścia przez wszystkie sprzedane bilety.
PostgreSQL z~indeksem robi to w~\val{postgres}{10000000}{on}{R6}\,ms (Index
Only Scan), MariaDB w~\val{mariadb}{10000000}{on}{R6}\,ms, MongoDB
w~\val{mongo}{10000000}{on}{R6}\,ms. Valkey odczytuje gotowy licznik ze zbioru
posortowanego w~\val{valkey}{10000000}{on}{R6}\,ms. Utrzymuje go przy każdej
sprzedaży, czyli przenosi koszt agregacji na zapis.

\paragraph{Valkey bez struktur indeksujących.} Zapytania po atrybucie
(historia klienta R4, wygaszanie U6, usuwanie D3, D4, D6, agregacja J2)
wymagają w~Valkey przejścia \texttt{SCAN} przez wszystkie klucze rezerwacji
i~zdekodowania każdego JSON-a. Przy 10 mln biletów to około
\val{valkey}{10000000}{off}{R4}\,ms na operację, najgorszy wynik w~całym badaniu.
Po dodaniu zbiorów \texttt{ix:*} ta sama historia klienta trwa
\val{valkey}{10000000}{on}{R4}\,ms. W~modelu klucz--wartość brak ,,indeksu''
nie oznacza wolniejszego planu, tylko brak jakiejkolwiek ścieżki dostępu.

\paragraph{Zmienność.} Mediany, odchylenia standardowe i~p95 są w~pliku
\texttt{generated/data/scenarios\_summary.csv}. Najwyższy współczynnik
zmienności mają operacje punktowe trwające ułamek milisekundy. Pojedyncze
powtórzenie trafiające na zimną pamięć podręczną albo punkt kontrolny WAL
wydłuża się wtedy kilkukrotnie. Dla operacji trwających powyżej 10\,ms
rozrzut jest mały i~nie zmienia kolejności baz.


\subsection{Wpływ indeksów}

\begin{figure}[H]
\centering
\begin{tikzpicture}
\begin{axis}[tb, ybar, bar width=1.6pt, ymode=log, log origin=infty, height=6.5cm,
  ylabel={przyspieszenie (bez / z~indeksami)}, xtick={0,...,26},
  xticklabels={C1,C2,C3,C4,C5,C6,R1,R2,R3,R4,R5,R6,U1,U2,U3,U4,U5,U6,D1,D2,D3,D4,D5,D6,J1,J2,J3},
  x tick label style={font=\scriptsize, rotate=90}, enlarge x limits=0.02,
  extra y ticks={1}, extra y tick style={grid=major, grid style={black, thick}}]
  \addplot[fill=cpg, draw=none] table[x=i, y=postgres] {generated/plots/speedup_10000000.dat}; \addlegendentry{PostgreSQL}
  \addplot[fill=cmaria, draw=none] table[x=i, y=mariadb] {generated/plots/speedup_10000000.dat}; \addlegendentry{MariaDB}
  \addplot[fill=cmongo, draw=none] table[x=i, y=mongo] {generated/plots/speedup_10000000.dat}; \addlegendentry{MongoDB}
  \addplot[fill=cvalkey, draw=none] table[x=i, y=valkey] {generated/plots/speedup_10000000.dat}; \addlegendentry{Valkey}
\end{axis}
\end{tikzpicture}
\caption{Ile razy scenariusz jest szybszy z~indeksami, 10 mln biletów (skala
logarytmiczna, wartość poniżej 1 oznacza spowolnienie). Brak słupka: brak
pomiaru bez indeksów (np. przekroczony limit czasu).}
\label{fig:speedup}
\end{figure}

Indeksy przyspieszają przede wszystkim scenariusze filtrujące po kolumnie
obcej. Przy 10 mln biletów historia klienta (R4) w~PostgreSQL skraca się
ze \val{postgres}{10000000}{off}{R4} do \val{postgres}{10000000}{on}{R4}\,ms,
w~MariaDB z~\val{mariadb}{10000000}{off}{R4} do \val{mariadb}{10000000}{on}{R4}\,ms,
a~w~MongoDB z~\val{mongo}{10000000}{off}{R4} do \val{mongo}{10000000}{on}{R4}\,ms.
Wygaszanie rezerwacji (U6) w~MariaDB spada z~\val{mariadb}{10000000}{off}{U6}
do \val{mariadb}{10000000}{on}{U6}\,ms.

Koszt indeksów widać przy wstawianiu. Import \num{10000} biletów (C3)
w~PostgreSQL wydłuża się z~\val{postgres}{10000000}{off}{C3} do
\val{postgres}{10000000}{on}{C3}\,ms, w~MariaDB z~\val{mariadb}{10000000}{off}{C3}
do \val{mariadb}{10000000}{on}{C3}\,ms, a~w~MongoDB z~\val{mongo}{10000000}{off}{C3}
do \val{mongo}{10000000}{on}{C3}\,ms. Każdy nowy wiersz trzeba wstawić także
do dwóch indeksów tabeli \texttt{bilety}. Pojedyncze zapisy po kluczu
(C1, U1, U2) zmieniają się w~granicach zmienności pomiaru. Indeksy
zwiększają też rozmiar bazy (tab.~\ref{tab:storage}) i~wymagają czasu na
budowę (tab.~\ref{tab:load}).

\subsection{Analiza planów zapytań}

Tabele~\ref{tab:plans-rel} i~\ref{tab:plans-nosql} zestawiają automatycznie
odczytane z~planów ścieżki dostępu (dla 13 scenariuszy odczytujących dane),
liczbę przeczytanych wierszy lub dokumentów i~czas wykonania raportowany przez
silnik. Komórki z~pełnym skanem są wyróżnione. Bez indeksów pełny skan
występuje w~\fullscans{postgres}{10000000}{off} z~\PlanScenarioCount{} planów
PostgreSQL, \fullscans{mariadb}{10000000}{off} MariaDB i~\fullscans{mongo}{10000000}{off}
MongoDB. Po dodaniu indeksów zostaje ich odpowiednio \fullscans{postgres}{10000000}{on},
\fullscans{mariadb}{10000000}{on} i~\fullscans{mongo}{10000000}{on}.

\begin{table}[htbp]
\centering
\scriptsize
\caption{Plany zapytań, bazy relacyjne, 10 mln biletów: ścieżki dostępu,
przeczytane wiersze, czas silnika. Tło: pełny skan tabeli.}
\label{tab:plans-rel}
\resizebox{\linewidth}{!}{\tbinput{generated/tables/plans_postgres_mariadb_10000000.tex}}
\end{table}

\begin{table}[htbp]
\centering
\scriptsize
\caption{Plany zapytań, MongoDB (\texttt{explain} z~\texttt{executionStats})
i~polecenia Valkey, 10 mln biletów.}
\label{tab:plans-nosql}
\resizebox{\linewidth}{!}{\tbinput{generated/tables/plans_mongo_valkey_10000000.tex}}
\end{table}

Wnioski z~planów:
\begin{itemize}
  \item W~R4 PostgreSQL po dodaniu indeksów zastępuje dwa skany sekwencyjne
  dwoma skanami indeksów (\texttt{rezerwacje\_klient},
  \texttt{bilety\_rezerwacja}). Czas silnika spada z~setek milisekund do
  dziesiątych części milisekundy (rys.~\ref{lst:r4}). MariaDB
  przechodzi z~\texttt{ALL} na \texttt{ref}: ponad 11 mln przeczytanych wierszy
  zmienia się w~kilka.
  \item R6 pokazuje, że indeks nie zawsze pomaga. PostgreSQL wykonuje
  Index Only Scan na \texttt{bilety(wydarzenie\_id, status)}, bez czytania
  tabeli. MariaDB czyta cały indeks (\texttt{index}). MongoDB nadal robi
  \texttt{COLLSCAN}, bo \texttt{\$match} po samym polu \texttt{status} nie
  pasuje do prefiksu indeksu złożonego \texttt{(wydarzenie\_id, status)}.
  Kolejność pól w~indeksie złożonym ma znaczenie.
  \item W~R2 z~indeksem PostgreSQL zostawia pełny skan małej tabeli
  \texttt{sektory} (kilkaset wierszy). Planer uznaje go za tańszy od
  indeksu, co jest poprawną decyzją kosztową, a~nie brakiem indeksu.
  \item W~J1 PostgreSQL używa indeksu GIN na \texttt{metadane} (Bitmap Index
  Scan), MongoDB indeksu na zagnieżdżonych polach \texttt{metadane.gatunek}
  i~\texttt{metadane.tagi}. Dane JSON są indeksowane na równi z~kolumnami.
  \item Valkey nie ma planów. Tabela pokazuje polecenia: bez struktur
  \texttt{ix:*} zostaje \texttt{SCAN} po przestrzeni kluczy, z~nimi
  \texttt{ZRANGE} po zbiorze posortowanym.
\end{itemize}

\begin{figure}[H]
{\VerbatimInput[firstline=9,frame=single,breaklines=true,breakanywhere=true,fontsize=\tiny,label={bez indeksów}]{data/results/plans/10000000/postgres/indexes-off/R4.txt}}
{\VerbatimInput[firstline=9,frame=single,breaklines=true,breakanywhere=true,fontsize=\tiny,label={z indeksami}]{data/results/plans/10000000/postgres/indexes-on/R4.txt}}
\caption{Plan R4 w~PostgreSQL bez indeksów i~z~indeksami, 10 mln biletów.
Surowe pliki z~\texttt{data/results/plans}, wstawione bez zmian.}
\label{lst:r4}
\end{figure}


\subsection{Narzut gRPC i śledzenia}
\label{sec:grpc}

\begin{table}[H]
\centering
\caption{Różnica między czasem wywołania mierzonym w~runnerze a~czasem
w~serwisie bazy, uśredniona po wszystkich scenariuszach i~rozmiarach.}
\label{tab:grpc}
\small
\tbinput{generated/tables/grpc.tex}
\end{table}

\begin{table}[H]
\centering
\caption{Narzut OpenTelemetry: mediana czasu wywołania (runner) dla tych
samych scenariuszy z~próbkowanym i~niepróbkowanym spanem głównym, 500 tys.
biletów, 30 powtórzeń.}
\label{tab:otel}
\small
\tbinput{generated/tables/overhead.tex}
\end{table}

Hop gRPC między runnerem a~serwisem bazy dodaje w~medianie dziesiąte części
milisekundy (tab.~\ref{tab:grpc}). Przy najszybszych operacjach, takich jak
odczyt rezerwacji po kluczu, to od jednej piątej do prawie połowy całego czasu
wywołania. Gdyby czas mierzył klient, różnice między bazami w~operacjach
punktowych byłyby rozmyte przez transport. Dlatego porównujemy czas mierzony
w~serwisie bazy. Średni narzut jest wyższy od mediany, bo obejmuje też
przesłanie dużych odpowiedzi (np. tysięcy wolnych miejsc w~R2).

Pomiar narzutu OpenTelemetry (tab.~\ref{tab:otel}) porównuje te same
scenariusze ze śladem próbkowanym i~niepróbkowanym, w~którym identyfikator
śladu nadal jest propagowany. Różnice mają oba znaki i~mieszczą się w~rozrzucie
pomiaru operacji trwających ułamek milisekundy. Pełne śledzenie z~eksportem
spanów nie zmienia więc wyników w~sposób mierzalny w~tym teście. Pozostaje
pomijalne także dlatego, że import hurtowy i~test współbieżności zapisują
spany tylko dla części operacji.

\subsection{Dane półstrukturalne (JSON)}
\label{sec:json}
%<*json>

Metadane wydarzeń, biletów, klientów, rezerwacji i~płatności są zapisane jako
JSON w~każdej bazie, choć na różne sposoby. PostgreSQL używa \texttt{JSONB}
(binarnie, z~operatorami i~indeksem GIN). MariaDB używa \texttt{JSON}
(tekst z~kontrolą poprawności, kolumna wirtualna do indeksu). W~MongoDB
metadane są natywnymi poddokumentami BSON. Valkey zapisuje JSON jako napis
w~polu hasza, a~częściową aktualizację wykonuje skrypt Lua z~biblioteką
\texttt{cjson}.

Filtr po polach JSON (J1) z~indeksem trwa przy 10 mln biletów
\val{postgres}{10000000}{on}{J1}\,ms w~PostgreSQL, \val{mariadb}{10000000}{on}{J1}\,ms
w~MariaDB, \val{mongo}{10000000}{on}{J1}\,ms w~MongoDB i~\val{valkey}{10000000}{on}{J1}\,ms
w~Valkey. Agregacja po polu \texttt{metadane.kanal} rezerwacji z~tygodnia (J2)
pokazuje koszt wyłuskiwania pola z~każdego dokumentu: PostgreSQL
\val{postgres}{10000000}{on}{J2}\,ms, MariaDB \val{mariadb}{10000000}{on}{J2}\,ms,
MongoDB \val{mongo}{10000000}{on}{J2}\,ms, Valkey \val{valkey}{10000000}{on}{J2}\,ms
(dekodowanie JSON-a po stronie klienta). Dopisanie elementu do tablicy w~JSON
(J3) jest operacją punktową we wszystkich bazach. PostgreSQL i~MariaDB
przepisują wtedy całą wartość kolumny, MongoDB wykonuje \texttt{\$push} na
miejscu, a~Valkey dekoduje i~koduje JSON w~skrypcie Lua.
%</json>

